% Root derivation, SRT1--24. All original measures and lattice labels retained.
\section{Shared rational periodization and its complete relative boundary}
\label{sec:shared-rational-theta}

\paragraph{Sources and scope.}
The adelic Schwartz source, rational periodization and original additive
measures are those of Alain Connes,
\href{https://arxiv.org/abs/math/9811068v1}{\emph{Trace formula in
noncommutative geometry and the zeros of the Riemann zeta function}},
Section III, and Alain Connes, Caterina Consani and Matilde Marcolli,
\href{https://arxiv.org/abs/math/0703392v1}{\emph{The Weil proof and the
geometry of the adeles class space}}, source labels \texttt{rangecVdef},
\texttt{varthetaH1V}, \texttt{fubini1}--\texttt{fubini4},
\texttt{nothappen} and \texttt{degmodifyV}. The latter paper explicitly
identifies the failure of absolute sum--integral interchange. We compute
that issue here for the actual GSB mixed antecedent and its entire
relative-rational transfer. The predecessor GSB and ABR proofs and the
following DMR proof accompany this section; no equivalence with a new
arithmetic cohomology is assumed. All tensors are algebraic and all
group-relation spaces are algebraic, without closure.

\subsection{Actual sources, convergence and the relative transfer}
Let $G=\mathbb Q^\times$ and retain
\begin{equation}\tag{SRT1}
\begin{gathered}
T=\mathcal S(\mathbb R)\otimes\bigotimes_p'
 (\mathcal S(\mathbb Q_p),\mathbf1_{\mathbb Z_p}),\qquad T_r=T^{\otimes r},\\
(U_a\xi)(z)=\xi(a^{-1}z),\qquad
E\xi(x)=\sum_{q\in G}\xi(qx),\quad x\in\mathbb A_\mathbb Q^\times,\\
E_\Delta t(x_1,\ldots,x_r)=
 \sum_{q\in G}t(qx_1,\ldots,qx_r),\qquad
E_r=E^{\otimes r}.
\end{gathered}
\end{equation}
Lebesgue measure at the real place and $\operatorname{vol}(\mathbb Z_p)=1$
at every finite place are used throughout. A pure elementary tensor and
a fixed idelic tuple impose $q\in D^{-1}\mathbb Z$ for some positive
integer $D$: the finite local supports have bounded denominators, and
outside a finite set the original reference indicators impose integrality.
Every nonzero real idelic coordinate has nonzero absolute value.
The real Schwartz bounds therefore give an absolutely convergent sum
over $q\ne0$. On compact subsets of the real coordinates bounded away
from zero the same bound applies to each real derivative, uniformly.
Finite sums of tensors retain these conclusions. Independent
periodization is absolutely convergent by the product of the individual
absolute sums. These arguments establish SRT1 on the original source,
not on an unspecified completion.
The algebraic tensor of the individual function images injects into
functions on the product set: for a finite linearly independent family
of functions one can select evaluation functionals spanning its dual
(otherwise a nonzero member would vanish at every point). Applying
those evaluations successively to each tensor factor proves injection.
Thus the function-space kernel of $E_r$ below is also its original
algebraic-tensor kernel.

Set $\mathcal V_\Delta=E_\Delta(T_r)$ as an actual space of functions
on the idelic tuples. There is an exact linear map on this image:
\begin{equation}\tag{SRT2}
\begin{gathered}
\mathcal P_r:\mathcal V_\Delta\longrightarrow E_r(T_r),\\
(\mathcal P_rF)(x_1,\ldots,x_r)=
 \sum_{(r_2,\ldots,r_r)\in G^{r-1}}
       F(x_1,r_2x_2,\ldots,r_rx_r),\qquad
\mathcal P_rE_\Delta=E_r.
\end{gathered}
\end{equation}
Indeed, the change of indices
$(q,r_2,\ldots,r_r)\mapsto(q,qr_2,\ldots,qr_r)$ is a bijection of
$G^r$. The double absolute sum is finite by the preceding estimate on
$E_r$, and the triangle inequality bounds the absolute sum of the
grouped values $F$ by it. This also proves independence from the chosen
antecedent of $F$. The map is onto its displayed codomain. Its exact
kernel is
\begin{equation}\tag{SRT3}
\ker\mathcal P_r=E_\Delta(\ker E_r),\qquad
\ker E_\Delta\subset\ker E_r.
\end{equation}
For if $F=E_\Delta t$, SRT2 gives $\mathcal P_rF=E_rt$ in both
directions. Every rational $r_i$ in SRT2 multiplies both the real and
all the finite coordinates. Omitting its finite coordinates would give
a different map.

For $r=2$ write $\mathcal R_pF(x,y)=F(x,py)$.
Common rational reindexing proves
\begin{equation}\tag{SRT4}
E_\Delta(U_p\otimes I)=\mathcal R_pE_\Delta,\qquad
E_\Delta(U_p\otimes U_p)=E_\Delta.
\end{equation}
The first identity follows by substituting $q=pq'$ in
$\sum_qt(qx/p,qy)$. The second uses the same substitution in both
arguments. These are identities of the original functions.

For comparison with the source's spherical coordinates, retain
\begin{equation}\tag{SRT5}
\begin{gathered}
\Theta(t_a\otimes h)(z)=h(z_\infty/a)
            \mathbf1_{a\widehat{\mathbb Z}}(z_f),\quad a\in\mathbb Q_{>0}^\times,
\quad h\in\mathcal S_{\rm even}(\mathbb R),\\
\ell=\prod_p p^{\max_i v_p(a_i)},\qquad x_i=(y_i,1_f),\quad y_i>0,\\
E_\Delta\left(\bigotimes_i\Theta(t_{a_i}\otimes h_i)\right)(x_i)
 =2\sum_{n\ge1}\prod_i h_i(\ell n y_i/a_i).
\end{gathered}
\end{equation}
The finite constraints are precisely $q\in\ell\mathbb Z$; the factor
two is the sum over the two rational signs. With
$M_+h(s)=\int_0^\infty h(u)u^{s-1}du$, Mellin integration gives
\begin{equation}\tag{SRT6}
\int_{(0,\infty)^r}E_\Delta(t)(y_i,1_f)
       \prod_i y_i^{s_i-1}dy_i
 =2\zeta\left(\sum_i s_i\right)
       \prod_i(a_i/\ell)^{s_i}M_+h_i(s_i).
\end{equation}
Here each individual absolute Mellin integral must converge and
$\operatorname{Re}\sum_i s_i>1$. In particular
$\operatorname{Re}s_i>0$ suffices for arbitrary even Schwartz $h_i$.
Absolute Fubini follows by replacing each $h_i$ by its absolute value
and using the convergent series $\sum_n n^{-\sum_i\operatorname{Re}s_i}$.
The original $\zeta$ and all finite rational scales remain in SRT6.

\subsection{The original mixed antecedent has a nonzero shared image}
Fix a rational prime $p$. Retain the exact GSB functions
\begin{equation}\tag{SRT7}
\begin{gathered}
\eta_0(u)=(1-2\pi u^2)e^{-\pi u^2},\qquad
\eta_1(u)=2\pi u^2e^{-\pi u^2},\\
\phi_{p,0}=(p\mathbf1_{p\mathbb Z_p}-\mathbf1_{\mathbb Z_p})/(p-1),
\quad
\phi_{p,1}=p(\mathbf1_{\mathbb Z_p}-\mathbf1_{p\mathbb Z_p})/(p-1),\\
a_p=\eta_1\otimes\phi_{p,0}\otimes
          \bigotimes_{l\ne p}\mathbf1_{\mathbb Z_l},\qquad
d_p=\eta_0\otimes\phi_{p,1}\otimes
          \bigotimes_{l\ne p}\mathbf1_{\mathbb Z_l},\\
z_p=\frac{((U_p-I)a_p)\otimes d_p}{p-1},\qquad
k_p=(U_p\otimes U_p-I)z_p.
\end{gathered}
\end{equation}
Rational reindexing gives $E_2z_p=0$. SRT4 shows
\begin{equation}\tag{SRT8}
E_\Delta z_p=(\mathcal R_p-I)E_\Delta(a_p\otimes d_p)/(p-1),
\quad \mathcal P_2 E_\Delta z_p=0,\quad E_\Delta k_p=0.
\end{equation}
Thus the shared image is a relative-rational difference. It is not
identified with the image of the coinvariant class $[k_p]$.

At $x_i=(y_i,1_f)$ the finite support of $d_p(qx_2)$ forces
$q\in\mathbb Z$ and $p\nmid q$. On that support
$\phi_{p,1}(q)=p/(p-1)$, $\phi_{p,0}(q)=-1/(p-1)$ and
$\phi_{p,0}(q/p)=0$. The original first summand in $z_p$ therefore
has zero contribution there, and its second summand gives
\begin{equation}\tag{SRT9}
F_p(y_1,y_2):=(E_\Delta z_p)(x_1,x_2)
 =\frac{2p}{(p-1)^3}\sum_{\substack{n\ge1\\p\nmid n}}
            \eta_1(ny_1)\eta_0(ny_2).
\end{equation}
For $y_1>0,y_2\ge1$ every term is strictly negative, so this is an
explicit nonzero element of $\ker\mathcal P_2$. The two real factors,
the two rational signs and every finite coefficient are retained.

Its full double Mellin formula is
\begin{equation}\tag{SRT10}
\begin{split}
\int_0^\infty\!\int_0^\infty F_p(y_1,y_2)y_1^{s-1}y_2^{t-1}dy_1dy_2
={}&\frac{2p}{(p-1)^3}(1-p^{-s-t})\zeta(s+t)\\
&\cdot\frac{2\pi\Gamma((s+2)/2)}{2\pi^{(s+2)/2}}\\
&\cdot\left\{\frac{\Gamma(t/2)}{2\pi^{t/2}}
       -\frac{2\pi\Gamma((t+2)/2)}{2\pi^{(t+2)/2}}\right\}.
\end{split}
\end{equation}
Absolute integration and termwise integration hold on the domain
$\operatorname{Re}s>-2$, $\operatorname{Re}t>0$,
$\operatorname{Re}(s+t)>1$: $\eta_1$ vanishes to order two at zero,
whereas $\eta_0(0)=1$. These conditions justify termwise integration
as in SRT6, now with the Euler factor deleting the multiples of $p$.
Direct Gaussian integration gives the two displayed Mellin factors.
The right side supplies a meromorphic continuation, retaining both
Gamma summands, their poles and the zero factor; no completed zeta
is substituted for the original function. At points of $s+t=1$ where
both displayed Gamma expressions are holomorphic, the coefficient of
$(s+t-1)^{-1}$, with $t$ fixed, is
\begin{equation}\tag{SRT11}
\frac{2p}{(p-1)^3}(1-p^{-1})
 \frac{2\pi\Gamma((s+2)/2)}{2\pi^{(s+2)/2}}
 \left\{\frac{\Gamma((1-s)/2)}{2\pi^{(1-s)/2}}
 -\frac{2\pi\Gamma((3-s)/2)}{2\pi^{(3-s)/2}}\right\}.
\end{equation}
For $0<s<1$ it is nonzero: writing $t=1-s\in(0,1)$, the
second term in braces is $t$ times the first by the Gamma recurrence.
This is the pole of this two-variable transform, not a claimed new
zero or pole of the original one-variable arithmetic quotient.
The formula also specifies the meromorphic residue coefficient on
the divisor. At intersections with Gamma poles it is not a claim
about the residue of a specialized one-variable function with a
possible higher-order pole.

\subsection{The exact diagonal moment and ray boundary}
Put $h_y(u)=\eta_1(y_1u)\eta_0(y_2u)$ and $A=y_1^2+y_2^2$.
The actual adelic diagonal pullback $\mu_y:T_2\to T$ gives
\begin{equation}\tag{SRT12}
(\mu_yz_p)(z)=\frac{p}{(p-1)^3}h_y(z_\infty)
   \mathbf1_{\widehat{\mathbb Z}}(z_f)
   \mathbf1_{\mathbb Z_p^\times}(\operatorname{pr}_p z).
\end{equation}
The finite additive integral contributes $1-p^{-1}$. On the real
line, Gaussian differentiation, with both original monomials, yields
\begin{equation}\tag{SRT13}
\begin{gathered}
I(y_1,y_2):=\int_0^\infty h_y(u)du
 =\frac{y_1^2}{2A^{3/2}}-\frac{3y_1^2y_2^2}{2A^{5/2}},\\
R_p(y):=\int_{\mathbb A_\mathbb Q}\mu_yz_p
 =\frac{2p}{(p-1)^3}(1-p^{-1})I(y_1,y_2)\\
 =\frac1{(p-1)^2}
       \left\{\frac{y_1^2}{A^{3/2}}
                      -\frac{3y_1^2y_2^2}{A^{5/2}}\right\}.
\end{gathered}
\end{equation}
In particular its sign is the sign of $y_1^2-2y_2^2$ for positive
$y_1,y_2$, and it vanishes precisely on that ray. Neither sign
nor vanishing determines the nonzero GSB class $[k_p]$.

Use the original Fourier transform
$\widehat h(\xi)=\int_\mathbb R h(u)e^{-2\pi i u\xi}du$.
Ordinary Poisson summation for the Schwartz function $h_y$, separately
at $\lambda$ and $p\lambda$, gives the exact identity
\begin{equation}\tag{SRT14}
\begin{split}
F_p(\lambda y_1,\lambda y_2)
=\frac{p}{(p-1)^3\lambda}\bigg\{&(1-p^{-1})\widehat h_y(0)
       +\sum_{m\ne0}\widehat h_y(m/\lambda)\\
       &-p^{-1}\sum_{m\ne0}\widehat h_y(m/(p\lambda))\bigg\}.
\end{split}
\end{equation}
Both complete sums are absolutely convergent; their zero summands are
displayed explicitly. Since $h_y$ is even and $h_y(0)=0$, SRT9 is
exactly the difference of the two integer theta sums multiplied by
$p/(p-1)^3$. Thus no unlisted endpoint term enters SRT14.
Schwartz bounds on $\widehat h_y$ prove, for every integer $N\ge0$,
\begin{equation}\tag{SRT15}
F_p(\lambda y_1,\lambda y_2)
 =\lambda^{-1}R_p(y)+O_{p,y,N}(\lambda^N)
 \quad(\lambda\downarrow0).
\end{equation}
The estimate follows by choosing a decay power exceeding $N+1$ in
the two nonzero Fourier sums. On compact subsets of $y_i>0$ those
Schwartz seminorms are uniform, including every specified finite
real derivative. At infinity the original Gaussian series has rapid
decay. The moment and the ray coefficient are therefore the same
explicit quantity, with its original additive and Fourier measures.

\subsection{Every relative rational and its exact cancellation}
Write a relative rational as $r=\pm m/n$, with $m,n\ge1$ coprime.
If $p\mid m$, the shared image at $(x_1,rx_2)$ is zero: the second
factor requires $v_p(q)=-v_p(m)<0$, outside both finite supports of
the first factor. Otherwise write $n=p^b n_0$, $p\nmid mn_0$.
All admissible summation rationals are exactly
$q=\pm p^b n_0l$ with $l\ge1$, $p\nmid l$. Define
\begin{equation}\tag{SRT16}
H_b(m,n_0;y)=2\sum_{\substack{l\ge1\\p\nmid l}}
       \eta_1(p^b n_0l y_1)\eta_0(ml y_2).
\end{equation}
Both signs of $q$ are in this factor two. Either sign of $r$
gives the same function since the original tests are even.
The finite coefficients are
\begin{equation}\tag{SRT17}
\begin{gathered}
\phi_{p,0}(q)=
  \begin{cases}-1/(p-1),&b=0,\\1,&b\ge1,\end{cases}
\qquad
\phi_{p,0}(q/p)=
  \begin{cases}0,&b=0,\\-1/(p-1),&b=1,\\1,&b\ge2,\end{cases}\\
\phi_{p,1}(qr)=p/(p-1),\qquad
F_b:=(E_\Delta z_p)(x_1,\pm m x_2/(p^bn_0))\\
=\begin{cases}
 pH_0/(p-1)^3,&b=0,\\
 \dfrac p{(p-1)^2}\{-H_0/(p-1)-H_1\},&b=1,\\
 \dfrac p{(p-1)^2}(H_{b-1}-H_b),&b\ge2.
\end{cases}
\end{gathered}
\end{equation}
These formulas follow by substituting the finite coefficients into
the two real summands of SRT7. In particular the exceptional $b=1$
term is not a geometric-series continuation of the $b=0$ term.
For each fixed $m,n_0$ and either rational sign, the finite sum is
\begin{equation}\tag{SRT18}
\sum_{b=0}^{B}F_b=-\frac p{(p-1)^2}H_B
 \quad(B\ge1),\qquad
\lim_{B\to\infty}H_B=0.
\end{equation}
The first two terms cancel their $H_0$ contributions; every following
term cancels the preceding $H_{b-1}$, leaving the displayed boundary.
The Gaussian factor in $\eta_1(p^Bn_0l y_1)$ proves the limit by an
absolutely summable bound. All reduced rationals with $p\nmid m$
are partitioned by these fixed coprime pairs $(m,n_0)$ and the two
signs. SRT2 proves absolute convergence over the entire rational set
at each fixed idelic pair, so this partition is legitimate. It
proves $\mathcal P_2 E_\Delta z_p=0$ with its finite-cutoff error
computed, as well as the independent reindexing proof in SRT8.

This pointwise absolute convergence does not justify interchanging
the complete relative sum with its ray boundary. For each $r$ let
$R_{p,r}(y)$ denote that individual ray limit. It exists by applying
SRT15 to the finitely many scaled terms in SRT17; it is zero when
$p\mid m$. For the subfamily
$r=1/n$, $p\nmid n$, SRT17 and SRT15 give
\begin{equation}\tag{SRT19}
\begin{split}
R_{p,1/n}(y)
&:=\lim_{\lambda\downarrow0}\lambda
     (E_\Delta z_p)(\lambda x_1,(1/n)\lambda x_2)\\
&=\frac{2}{(p-1)^2}I(ny_1,y_2)\\
&=\frac1{(p-1)^2}
 \left\{\frac{n^2y_1^2}{(n^2y_1^2+y_2^2)^{3/2}}
 -\frac{3n^2y_1^2y_2^2}{(n^2y_1^2+y_2^2)^{5/2}}\right\}\\
&=\frac1{(p-1)^2 n y_1}+O_{p,y}(n^{-3}).
\end{split}
\end{equation}
Here $\lambda x_i$ means multiplication at the real place only,
whereas $1/n$ acts at every place. The finite coordinates in SRT19
are therefore still those of SRT17. The last estimate follows by
Taylor's formula for $(1+u)^{-3/2}$ and $(1+u)^{-5/2}$ on a bounded
interval with $u=y_2^2/(n^2y_1^2)$; the full two summands remain above.
The leading term is positive. Consequently
\begin{equation}\tag{SRT20}
\sum_{r\in\mathbb Q^\times}|R_{p,r}(y)|=\infty,
\qquad
\lim_{\lambda\downarrow0}\lambda\,
       (\mathcal P_2E_\Delta z_p)(\lambda x_1,\lambda x_2)=0.
\end{equation}
Indeed $\sum_{p\nmid n}1/n$ diverges, whereas $\sum n^{-3}$
converges. Thus there is no absolutely convergent residue sum equal
to the second expression. One may instead take each fixed
$(m,n_0)$ chain first: its residue analogue of SRT18 has terminal
coefficient tending to zero as $B\to\infty$. This specifies an
ordered grouping; it is not an unconditional sum or an application
of absolute Fubini.

\subsection{The full original Euler factors in the comparison}
For $t=(h_1\mathbf1_{\widehat{\mathbb Z}})\otimes
(h_2\mathbf1_{\widehat{\mathbb Z}})$ with
$h_1,h_2\in\mathcal S_{\rm even}(\mathbb R)$, the complete finite support
calculation at $r=\pm m/n$ reduced gives
\begin{equation}\tag{SRT21}
E_\Delta t(x_1,\pm m x_2/n)
 =2\sum_{l\ge1}h_1(nl y_1)h_2(ml y_2).
\end{equation}
The two signs of the relative rational add another factor two.
In $\operatorname{Re}s>1$, $\operatorname{Re}t>1$, the arithmetic
coefficient is exactly
\begin{equation}\tag{SRT22}
\begin{gathered}
\mathcal C(s,t)=\sum_{\substack{m,n\ge1\\(m,n)=1}}n^{-s}m^{-t},\\
\mathcal C(s,t)=\prod_p\left\{1+
       \sum_{a\ge1}p^{-as}+\sum_{b\ge1}p^{-bt}\right\}
 =\prod_p\frac{1-p^{-s-t}}{(1-p^{-s})(1-p^{-t})},\\
\zeta(s+t)\mathcal C(s,t)=\zeta(s)\zeta(t).
\end{gathered}
\end{equation}
Absolute convergence justifies prime factorization; coprimality
means that at each prime at most one of the two exponents is positive.
Equivalently every ordered pair of positive integers is uniquely a
common positive factor times a coprime pair. This proves the last
identity directly without dividing by a possibly vanishing factor.
The shared Mellin value and the full relative transfer are therefore
\begin{equation}\tag{SRT23}
\begin{gathered}
2\zeta(s+t)M_+h_1(s)M_+h_2(t),\\
4\zeta(s+t)\mathcal C(s,t)M_+h_1(s)M_+h_2(t)
 =4\zeta(s)\zeta(t)M_+h_1(s)M_+h_2(t).
\end{gathered}
\end{equation}
These are identities on the absolute-convergence domain, followed
where appropriate by meromorphic continuation. At a meromorphic
point the divisor identity keeps the sum of all orders on the two
sides; cancellation of a zero and pole is not permission to erase
their individual factors. The relative transfer has nonzero kernel
by SRT9, so neither SRT22 nor a scalar ratio on a nonvanishing locus
is an inverse to that function-space map.

\subsection{Full support and the exact arithmetic scope}
For the original bounded distributive lattice $L$ retain
\begin{equation}\tag{SRT24}
\begin{gathered}
G_L(M)=\{(0,\lambda):\lambda\in L\}\cup(M\times\{1_L\}),
\quad \tau=(0,0_L),\quad e=(0,1_L),\\
A^\uparrow(v,\lambda)=(Av,\lambda).
\end{gathered}
\end{equation}
Addition has label join and scalar multiplication has label meet,
as in the originating
\href{https://github.com/KokunoYumeto/zeta-function-research-reader/blob/a2851f9eb352e7e4376bc629de86fa4ed9c1c434/workbenches/splitzero-tandem/foundations/split-support-geometry-v11/split_support_geometry_arithmetic_curve_v11.tex}{\emph{Split Support Geometry}, v11,
Definition \texttt{def:lattice-split}, by The Clankers}.
Every linear map above has the exact common-label lift SRT24.
Linearity proves compatibility with scalar amplitudes and joins;
zero scalars retain the original meet label on both sides. Thus
SRT2 becomes an equality of $G_L$-maps, SRT3 becomes its
same-label kernel statement, and the nonzero shared image of $z_p$
maps to $e$ at top support rather than to $\tau$. Every lower
zero label remains unchanged. The universal tensor and its
common-label comparison are the proved GFT maps; two independently
chosen labels are not silently identified with one common label.
Analytic convergence is proved in the original amplitude function
spaces before lifting the resulting linear maps. No infinite join
in an arbitrary lattice $L$ is postulated.

The exact diagonal pullback and its moment are the DMR2 map, with
the moment defined on the entire original tensor space. DMR16--23
calculates its descent on the joint-kernel coinvariants: changing
a GSB antecedent by an element of that kernel can alter the raw
prime integral coordinate arbitrarily, while the refined source
retains the entire coordinate separately. Its canonical ABR
receiver is zero on the whole GSB image. Accordingly the nonzero
shared function, its sign-changing ray residue and the divergent
absolute residue sum proved here are exact information about the
source and transfer; they are not a negative value of the original
Weil pairing. No positivity, purity or RH conclusion has been
inserted as an assumption.
